\section{Method}
\label{sec:method}

Our compiler loop has four stages: (i)~import asset provenance and link structure; (ii)~generate or import primitive candidates from heuristics, native lanes, or decomposition outputs; (iii)~run Newton diagnostics; (iv)~accept, reject, or fallback with an auditable reason.
The loop deliberately separates \emph{candidate generation} from \emph{simulation acceptance}.
This separation is the central engineering choice: geometry methods provide useful packages, but the compiler only promotes a package after named probes accept it under recorded settings.

\paragraph{Candidate generation.}
Each asset is first normalized into a package manifest that preserves source asset identity, stage path, units, and, when present, USD link boundaries~\cite{openusd2026}.
Rigid assets can receive bounding primitives, native primitive lanes, CPD-style primitive packages~\cite{cpd2026}, and convex-mesh baselines such as CoACD and V-HACD~\cite{coacd2022,vhacd}.
For articulated assets, candidate generation is link-framed: a primitive is emitted under the link that owns its source geometry, and the package records whether any primitive crosses a joint boundary.
Candidate quality is summarized by primitive or hull count, geometric residual when available, and editability, but these signals do not imply simulation acceptance.

\paragraph{Acceptance layer.}
The acceptance layer converts each candidate package into a Newton task scene and records a pass/fail label with configuration, seed, and package metadata~\cite{newton2026}.
Body-state and drop/settle gates label residual speed and compound COM/inertia risk.
Contact probes observe false clearance, blocked openings, excessive contact counts, and stack-or-slide behavior.
For robots, a link-boundary audit forbids cross-link merges; articulation smoke gates joint import, gravity hold, and simple trajectories before any broader robot claim.
The current generated-package robot task smoke verifies package consumption in that bounded task path.
It disables generated self-collision pairs for the smoke and does not establish whole-robot collision quality, manipulation performance, or end-effector accuracy.

\paragraph{Fallback policy.}
Failed packages route to another candidate lane, a convex-mesh fallback, SDF or manual review depending on configuration.
Fallback is recorded as an explicit outcome, not hidden as an omitted result.
An opt-in package body-state guard can fallback only flagged packages while preserving unflagged packages in the same task smoke, as recorded for capped bed versus Franka cylinders.
This keeps the selector conservative: accepted primitive packages remain native and inspectable, while rejected packages retain a reproducible reason.

\paragraph{Evidence registry.}
Quantitative statements in this manuscript must match the shared \texttt{claims.yaml} registry and dated records under \texttt{docs/records/}.
If a table cell has no dated record, the paper treats it as an open item rather than an implied result.
